% pdf-header.tex
% Extra styling for Pandoc-generated PDFs
% (included via --include-in-header)

\makeatletter
\@ifundefined{savenotes}{}{%%
  \let\savenotes\relax%%
  \let\endsavenotes\relax%%
}
\makeatother

% Hyphenation and line-breaking controls
% Reduce excessive hyphenation for better readability
\hyphenpenalty=500       % Default is 50 (lower = more hyphenation)
\exhyphenpenalty=500     % Penalty for breaking already-hyphenated words
\tolerance=1000          % Default is 200 (higher = more tolerance for loose spacing)
\emergencystretch=3em    % Allow extra stretch to avoid overfull boxes

\usepackage{needspace}
\usepackage{fontspec}
\usepackage{newunicodechar}
\usepackage{fancyvrb}
\usepackage{fvextra} % allow breaklines for Pandoc's Verbatim blocks
\usepackage{framed}
\usepackage{enumitem}
\usepackage{ltablex}  % Extends longtable with tabularx features
\fvset{breaklines=true,breakanywhere=true}
% Alias the standard \begin{verbatim} environment to fancyvrb's \begin{Verbatim}
% so that it also inherits breaklines=true from \fvset above.  Without this,
% indented code blocks (which Pandoc renders as \begin{verbatim}) silently
% overflow the page margin on long single-line content.
\RecustomVerbatimEnvironment{verbatim}{Verbatim}{}
\csname tracinglostchars\endcsname=1

% --------------------------------------------------
% Force all longtable columns to wrap text properly
% Redefine column types to use ragged-right paragraph columns
% --------------------------------------------------
\usepackage{array}
\renewcommand{\arraystretch}{1.2}
% Make 'l' columns behave as wrapping paragraph columns in longtable
\newcolumntype{L}[1]{>{\raggedright\arraybackslash}p{#1}}
% Default table font slightly smaller for better fit
\AtBeginEnvironment{longtable}{\small}

\providecommand{\BodyFontFamily}{TeX Gyre Pagella}
\providecommand{\HeadingFontFamily}{TeX Gyre Pagella}
\providecommand{\MonoFontFamily}{DejaVu Sans Mono}
\providecommand{\MonoFontFallbackA}{Inconsolata}
\providecommand{\MonoFontFallbackB}{TeX Gyre Cursor}
\providecommand{\ScreenGeometryOptions}{paperwidth=6in, paperheight=9in, margin=1.1in, twoside=false}
\providecommand{\PrintGeometryOptions}{paperwidth=6in, paperheight=9in, left=0.75in, right=0.375in, top=0.5in, bottom=0.5in, twoside=false}

% Provide fallback Highlighting/Shaded environments if Pandoc's template omits them
\makeatletter
\@ifundefined{VerbBar}{\newcommand{\VerbBar}{|}}{}
\@ifundefined{VERB}{\newcommand{\VERB}{\Verb[commandchars=\\\{\}]}}{}
\@ifundefined{Highlighting}{%
  \DefineVerbatimEnvironment{Highlighting}{Verbatim}{commandchars=\\\{\}}
}{}
\@ifundefined{Shaded}{%
  \newenvironment{Shaded}{\begin{snugshade}}{\end{snugshade}}
}{}
 \@ifundefined{abstract}{%
   \newenvironment{abstract}{%
     \begin{quote}
     \begingroup\centering\bfseries Abstract\par\endgroup\smallskip
   }{%
     \end{quote}
   }
 }{}
\makeatother

% Reconfigure Pandoc's Highlighting environment so literal includes wrap lines
\RecustomVerbatimEnvironment{Highlighting}{Verbatim}{%
  commandchars=\\\{\},% keep Pandoc's command chars intact
  breaklines=true,% wrap long lines instead of overflowing the margin
  breakanywhere=true
}
\defaultfontfeatures{Ligatures=TeX,Scale=MatchLowercase,Renderer=Harfbuzz}

% --------------------------------------------------
% TOC list spacing (PDF rendering)
% --------------------------------------------------
\newenvironment{tocblock}{%
  \begingroup%
  \setlist[itemize,1]{label={},labelsep=0pt,leftmargin=2.1em,itemindent=0pt,listparindent=0pt,align=left,itemsep=0pt,topsep=0.2em,parsep=0pt}%
  \setlist[itemize,2]{label=\textbullet,labelsep=0.35em,labelwidth=1em,leftmargin=2.6em,itemindent=0pt,align=left,itemsep=0pt,topsep=0.2em,parsep=0pt}%
  \setlist[itemize,3]{label=\textbullet,labelsep=0.35em,labelwidth=1em,leftmargin=3.4em,itemindent=0pt,align=left,itemsep=0pt,topsep=0.2em,parsep=0pt}%
  \setlist[enumerate,1]{labelsep=0.35em,labelwidth=2.1em,leftmargin=2.6em,align=right,itemindent=0pt,itemsep=0pt,topsep=0.2em,parsep=0pt}%
  \setlist[enumerate,2]{labelsep=0.35em,labelwidth=1.6em,leftmargin=2.9em,align=right,itemindent=0pt,itemsep=0pt,topsep=0.2em,parsep=0pt}%
  \setlist[enumerate,3]{labelsep=0.35em,labelwidth=1.6em,leftmargin=3.7em,align=right,itemindent=0pt,itemsep=0pt,topsep=0.2em,parsep=0pt}%
}{\endgroup}

% --------------------------------------------------
% PDF mode switching: 'print' vs 'screen'
% --------------------------------------------------
\usepackage{ifthen}
\usepackage{etoolbox}

\providecommand{\RenderGrayscaleScreen}{0}
\providecommand{\RenderGrayscalePrint}{0}
\providecommand{\BodyRunningHeadersEnabled}{0}
\providecommand{\AppendixRunningHeadersEnabled}{0}
\providecommand{\RunningHeaderScope}{chapter}
\providecommand{\AppendixRunningHeaderScope}{chapter}

% --------------------------------------------------
% Font configuration (XeLaTeX)
% --------------------------------------------------
\newif\ifiaspMonoFontChosen
\iaspMonoFontChosenfalse

\newcommand{\iaspTrySetMonoFont}[1]{%%
  \ifiaspMonoFontChosen
  \else
    \IfFontExistsTF{#1}{%%
      \setmonofont{#1}%%
      \iaspMonoFontChosentrue%%
    }{}
  \fi
}

\newcommand{\iaspSetMonospacedFont}{%%
  \iaspMonoFontChosenfalse%%
  \iaspTrySetMonoFont{\MonoFontFamily}%%
  \ifiaspMonoFontChosen
  \else
    \PackageWarningNoLine{iasp-fonts}{Mono font '\MonoFontFamily' missing; trying fallbacks.}%%
    \iaspTrySetMonoFont{\MonoFontFallbackA}%%
    \iaspTrySetMonoFont{\MonoFontFallbackB}%%
  \fi
  \ifiaspMonoFontChosen
  \else
    \PackageWarningNoLine{iasp-fonts}{No configured mono font found; falling back to 'Latin Modern Mono'.}%%
    \setmonofont{Latin Modern Mono}%%
  \fi
}

\setmainfont{\BodyFontFamily}
\newfontfamily\HeadingFont{\HeadingFontFamily}
\iaspSetMonospacedFont


\newcommand{\setnoteglyphfont}{%
  \IfFontExistsTF{Noto Color Emoji}{%
    \newfontfamily\noteglyphfont{Noto Color Emoji}[Renderer=Harfbuzz,Scale=MatchLowercase]%
  }{%
    \IfFontExistsTF{Noto Emoji}{%
      \newfontfamily\noteglyphfont{Noto Emoji}[Renderer=Harfbuzz,Scale=MatchLowercase]%
    }{%
      \IfFontExistsTF{Symbola}{%
        \newfontfamily\noteglyphfont{Symbola}[Scale=MatchLowercase]%
      }{%
        \newfontfamily\noteglyphfont{DejaVu Sans}[Scale=MatchLowercase]%
      }%
    }%
  }%
}
\setnoteglyphfont

\newcommand{\setsymbolfont}{%
  \IfFontExistsTF{Noto Sans Symbols2}{%
    \newfontfamily\symbolfont{Noto Sans Symbols2}[Scale=MatchLowercase]%
  }{%
    \IfFontExistsTF{Noto Sans Symbols}{%
      \newfontfamily\symbolfont{Noto Sans Symbols}[Scale=MatchLowercase]%
    }{%
      \IfFontExistsTF{DejaVu Sans}{%
        \newfontfamily\symbolfont{DejaVu Sans}[Scale=MatchLowercase]%
      }{%
        \newfontfamily\symbolfont{Symbola}[Scale=MatchLowercase]%
      }%
    }%
  }%
}
\setsymbolfont

% Ensure U+25D0 (◐) renders from a symbol-capable font
\newunicodechar{◐}{{\symbolfont◐}}

\newcommand{\setcheckboxfont}{%
  \IfFontExistsTF{Libertinus Serif}{%
    \newfontfamily\checkboxfont{Libertinus Serif}[Scale=MatchLowercase]%
  }{%
    \IfFontExistsTF{Libertinus Serif Regular}{%
      \newfontfamily\checkboxfont{Libertinus Serif Regular}[Scale=MatchLowercase]%
    }{%
      \IfFontExistsTF{STIX Two Text}{%
        \newfontfamily\checkboxfont{STIX Two Text}[Scale=MatchLowercase]%
      }{%
        \newfontfamily\checkboxfont{DejaVu Sans}[Scale=MatchLowercase]%
      }%
    }%
  }%
}
\setcheckboxfont

\newcommand{\emojiFallback}[1]{{\noteglyphfont\symbol{#1}}}
\newunicodechar{📺}{\emojiFallback{"1F4FA}}
\newunicodechar{🌐}{\emojiFallback{"1F310}}
\newunicodechar{📘}{\emojiFallback{"1F4D8}}
\newunicodechar{⭐}{\emojiFallback{"2B50}}
\newunicodechar{🗒}{\emojiFallback{"1F5D2}}   % spiral notepad
\newunicodechar{📊}{\emojiFallback{"1F4CA}}   % bar chart
\newunicodechar{🐣}{\emojiFallback{"1F423}}   % hatching chick
\newunicodechar{☐}{{\checkboxfont\symbol{"2610}}}
\newunicodechar{☑}{{\checkboxfont\symbol{"2611}}}

% Note variant glyphs - fallback to noteglyphfont when encountered in prose
\newunicodechar{💡}{\emojiFallback{"1F4A1}}   % idea
\newunicodechar{🔍}{\emojiFallback{"1F50D}}   % diagnostic
\newunicodechar{⇥}{{\noteglyphfont ⇥}}        % transition
\newunicodechar{⚙}{{\noteglyphfont ⚙}}        % usage
\newunicodechar{⚖}{{\noteglyphfont ⚖}}        % responsibility
\newunicodechar{⟲}{{\noteglyphfont ⟲}}        % reflexive
\newunicodechar{⇢}{{\noteglyphfont ⇢}}        % navigation
\newunicodechar{✕}{{\noteglyphfont ✕}}        % misuse
\newunicodechar{◉}{{\noteglyphfont ◉}}        % context
\newunicodechar{◎}{{\noteglyphfont ◎}}        % scope
\newunicodechar{⊢}{{\noteglyphfont ⊢}}        % methodology
\newunicodechar{⊕}{{\noteglyphfont ⊕}}        % synthesis
\newunicodechar{◊}{{\noteglyphfont ◊}}        % disclosure
\newunicodechar{⇆}{{\noteglyphfont ⇆}}        % contrast
\newunicodechar{∎}{{\noteglyphfont ∎}}        % closing
\newunicodechar{⟷}{{\noteglyphfont ⟷}}        % bridge
\newunicodechar{↗}{{\noteglyphfont ↗}}        % breadcrumb
\newunicodechar{⁂}{{\noteglyphfont ⁂}}        % insight
\newunicodechar{※}{{\noteglyphfont ※}}        % warning
\newunicodechar{∴}{{\noteglyphfont ∴}}        % note (default)
\newunicodechar{∞}{{\noteglyphfont ∞}}        % integration

\newcommand{\sectionheadingglyph}[1]{%%
  \llap{{\noteglyphfont #1}\hspace{0.6em}}%%
}

\newcommand{\collectionheadingglyph}[1]{%%
  \llap{{\noteglyphfont #1}\hspace{0.35em}}%% tighter hang for collection-level glyphs
}

% In print mode, glyphs must NOT hang into the gutter (KDP margin enforcement).
% Redefine after isprint is known; \AtBeginDocument ensures it overrides any
% later redefinitions from included packages.
\newcommand{\iaspRedefineGlyphMacrosForPrint}{%%
  \renewcommand{\sectionheadingglyph}[1]{%%
    {\noteglyphfont ##1}\hspace{0.4em}%%
  }%%
  \renewcommand{\collectionheadingglyph}[1]{%%
    {\noteglyphfont ##1}\hspace{0.3em}%%
  }%%
}

% Pandoc will set \pdfmode via --variable=pdfmode:MODE or metadata.
% We interpret that into booleans for use below. If \pdfmode is
% undefined or not 'print', we default to 'screen'.

\newboolean{isprint}
\newboolean{isscreen}
\newif\ifiaspinappendixmode
\iaspinappendixmodefalse
\newif\ifiaspChapterMarkSeen
\iaspChapterMarkSeenfalse

% default to screen
\setboolean{isprint}{false}
\setboolean{isscreen}{true}

% If the user specified pdfmode=print, flip the flags.
\ifdefined\pdfmode
  \ifthenelse{\equal{\pdfmode}{print}}{%%
    \setboolean{isprint}{true}%%
    \setboolean{isscreen}{false}%%
  }{}
\fi

\newboolean{rendergrayscale}
\setboolean{rendergrayscale}{false}

\ifthenelse{\boolean{isprint}}{%%
  \ifthenelse{\equal{\RenderGrayscalePrint}{1}}{%%
    \setboolean{rendergrayscale}{true}%%
  }{}
}{%%
  \ifthenelse{\equal{\RenderGrayscaleScreen}{1}}{%%
    \setboolean{rendergrayscale}{true}%%
  }{}
}

% --------------------------------------------------
% Font size reduction for KDP print builds
% --------------------------------------------------
% Apply 9pt font in print mode to reduce pagination for KDP's 828-page limit
% Use \AtBeginDocument to ensure it's applied after all packages have loaded
\ifthenelse{\boolean{isprint}}{%%
  \AtBeginDocument{%%
    \renewcommand{\normalsize}{\fontsize{9pt}{11pt}\selectfont}%%
    \renewcommand{\small}{\fontsize{8pt}{9.5pt}\selectfont}%%
    \renewcommand{\footnotesize}{\fontsize{7.5pt}{9pt}\selectfont}%%
    \renewcommand{\large}{\fontsize{10.5pt}{12.5pt}\selectfont}%%
    \renewcommand{\Large}{\fontsize{12pt}{14pt}\selectfont}%%
    \renewcommand{\LARGE}{\fontsize{14pt}{17pt}\selectfont}%%
    \renewcommand{\huge}{\fontsize{17pt}{20pt}\selectfont}%%
    \renewcommand{\Huge}{\fontsize{20pt}{24pt}\selectfont}%%
    \normalsize  %% Apply immediately
    \iaspRedefineGlyphMacrosForPrint  %% Inline glyphs: no llap into gutter
  }%%
}{}



% --------------------------------------------------
% Page layout
% --------------------------------------------------
% Note: geometry package is loaded by Pandoc's template with options
% passed via --variable=geometry: in build_book.py
% The \PrintGeometryOptions and \ScreenGeometryOptions are kept for reference
% but no longer used directly
\usepackage{microtype} % nicer text spacing
\usepackage{xurl}      % allow URL line-breaks at slashes, dots, etc.
\urlstyle{same}        % keep URL font consistent with surrounding text
\muskipdef\Urlmuskip=0
\Urlmuskip=0mu plus 2mu\relax % permit additional URL break flexibility
\makeatletter
% Add extra safe breakpoints for long query-heavy URLs in references.
\g@addto@macro\UrlBreaks{\do\-\do\_\do\~\do\&\do\=\do\?\do\#}
\makeatother

% --------------------------------------------------
% Page style (headers/footers)
% --------------------------------------------------
\usepackage{fancyhdr}
\newcommand{\defaultpagestylename}{fancy}

\newcommand{\iaspRunningHeadFormat}[1]{\small\itshape\nouppercase{#1}}

\newcommand{\iaspPartHeaderText}{}
\newif\ifiaspParthasvalue
\iaspParthasvaluefalse

\newcommand{\iaspSetPartHeader}[1]{%%
  \gdef\iaspPartHeaderText{#1}%%
  \iaspParthasvaluetrue%%
}

\newcommand{\iaspClearPartHeader}{%%
  \gdef\iaspPartHeaderText{}%%
  \iaspParthasvaluefalse%%
}

\newif\ifiaspNextPartHeaderPending
\iaspNextPartHeaderPendingfalse
\newcommand{\iaspNextPartHeaderText}{}

\newcommand{\iaspSetNextPartHeader}[1]{%%
  \gdef\iaspNextPartHeaderText{#1}%%
  \iaspNextPartHeaderPendingtrue%%
}

\newcommand{\iaspApplyNextPartHeader}{%%
  \ifiaspNextPartHeaderPending
    \iaspSetPartHeader{\iaspNextPartHeaderText}%%
    \iaspNextPartHeaderPendingfalse%%
  \fi
}

\newcommand{\iaspRunningHeadSource}{%%
  \ifthenelse{\equal{\RunningHeaderScope}{part}}{%%
    \ifiaspParthasvalue
      \iaspPartHeaderText%%
    \else
      \leftmark%%
    \fi
  }{%%
    \leftmark%%
  }%%
}

\newcommand{\iaspAppendixRunningHeadSource}{%%
  \ifthenelse{\equal{\AppendixRunningHeaderScope}{section}}{%%
    \leftmark%%
  }{%%
    \ifiaspChapterMarkSeen
      \iaspLastChapterMark%%
    \else
      \leftmark%%
    \fi
  }%%
}

\newcommand{\iaspRunningHeadText}{%%
  \ifiaspinappendixmode
    \ifnum\AppendixRunningHeadersEnabled=1 \iaspRunningHeadFormat{\iaspAppendixRunningHeadSource}\fi%%
  \else
    \ifnum\BodyRunningHeadersEnabled=1 \iaspRunningHeadFormat{\iaspRunningHeadSource}\fi%%
  \fi
}

% For print, folio is at the outer margin edge; running head at the inner edge.
%   Even (verso): folio outer-left [LE], running head inner-right [RE]
%   Odd  (recto): running head inner-left [LO], folio outer-right [RO]
% They occupy separate header zones with the full column width between them.
% Footer is empty; no wasted space at the bottom.
% For screen, keep things subdued: footer page number only,
% and no horizontal head rule.
\ifthenelse{\boolean{isprint}}{%%
  \pagestyle{fancy}
  \fancyhf{}% clear all
  % Verso: folio at outer-left, running head at inner-right
  \fancyhead[LE]{\small\thepage}
  \fancyhead[RE]{\iaspRunningHeadText}
  % Recto: running head at inner-left, folio at outer-right
  \fancyhead[LO]{\iaspRunningHeadText}
  \fancyhead[RO]{\small\thepage}
  % Keep default \headrulewidth (visible horizontal line)
}{%%
  % Screen: minimal; running heads are optional and disabled by default
  \pagestyle{fancy}
  \fancyhf{}
  \fancyfoot[C]{\thepage}
  \fancyhead[LE,RO]{\iaspRunningHeadText}
  \renewcommand{\headrulewidth}{0pt}
}

\fancypagestyle{plain}{%% ensure chapter-opening pages follow the same rules
  \fancyhf{}
  \ifthenelse{\boolean{isprint}}{%%
    \fancyhead[LE]{\small\thepage}%%
    \fancyhead[RE]{\iaspRunningHeadText}%%
    \fancyhead[LO]{\iaspRunningHeadText}%%
    \fancyhead[RO]{\small\thepage}%%
    \renewcommand{\headrulewidth}{0.4pt}%%
  }{%%
    \fancyfoot[C]{\thepage}%%
    \fancyhead[LE,RO]{\iaspRunningHeadText}%%
    \renewcommand{\headrulewidth}{0pt}%%
  }
}


% Make sure \chaptermark exists even for 'article' class
\providecommand{\chaptermark}[1]{}

\newcommand{\iaspLastChapterMark}{}
\newcommand{\iaspLastSectionMark}{}

\newif\ifiaspChapterMarkOverridePending
\iaspChapterMarkOverridePendingfalse
\newcommand{\iaspChapterMarkOverrideText}{}

\newcommand{\iaspSetNextChapterRunningHead}[1]{%%
  \gdef\iaspChapterMarkOverrideText{#1}%%
  \iaspChapterMarkOverridePendingtrue%%
}

\newcommand{\iaspSetRunningHeadNow}[1]{%%
  \iaspSetPartHeader{#1}%%
  \gdef\iaspLastChapterMark{\unexpanded{#1}}%%
  \markboth{\unexpanded{#1}}{}%%
  \iaspChapterMarkSeentrue%%
}

\newcommand{\iaspClearRunningHead}{%%
  \gdef\iaspLastChapterMark{}%%
  \markboth{}{}%%
}

% Use chapter titles for \leftmark (no section marks in header)
\renewcommand{\chaptermark}[1]{%%
  \iaspApplyNextPartHeader%% ensure deferred part headers activate after the prior page flushes
  \ifiaspChapterMarkOverridePending%%
    \gdef\iaspLastChapterMark{\unexpanded{\iaspChapterMarkOverrideText}}%%
    \markboth{\unexpanded{\iaspChapterMarkOverrideText}}{}%%
    \iaspChapterMarkOverridePendingfalse%%
  \else%%
    \gdef\iaspLastChapterMark{\unexpanded{#1}}%%
    \markboth{\unexpanded{#1}}{}%%
  \fi%%
  \iaspChapterMarkSeentrue%%
}
\renewcommand{\sectionmark}[1]{%%
  \gdef\iaspLastSectionMark{#1}%%
  \markboth{#1}{}%%
}

% Make sure header area is tall enough
\setlength{\headheight}{14pt}
\setlength{\headsep}{12pt}

\newcommand{\includeBannerText}{}
\newif\ifincludeBannerActive
\includeBannerActivefalse

\newcommand{\includeBannerStripe}{%%
  \begingroup
  \setlength{\fboxsep}{6pt}%%
  \colorbox{includebannerbg}{%%
    \parbox{\dimexpr\textwidth-2\fboxsep\relax}{\centering\footnotesize\sffamily\textbf{\includeBannerText}}%%
  }%%
  \endgroup
}

\fancypagestyle{include-banner}{%%
  \fancyhf{}
  \fancyhead[C]{\includeBannerStripe}
  \fancyfoot[C]{\thepage}
  \renewcommand{\headrulewidth}{0pt}
}

\newcommand{\activateIncludeBanner}[1]{%%
  \def\includeBannerText{#1}%%
  \includeBannerActivetrue%%
  \pagestyle{include-banner}%%
}

\newcommand{\deactivateIncludeBanner}{%%
  \includeBannerActivefalse%%
  \pagestyle{\defaultpagestylename}%%
  \def\includeBannerText{}%%
}

\newenvironment{includepagebanner}[1]{%%
  \clearpage%% ensure bannered inserts always start on a fresh page
  \activateIncludeBanner{#1}%%
    \thispagestyle{include-banner}%%
}{%%
  \deactivateIncludeBanner%%
}

\newenvironment{includepagebannerchapter}[1]{%%
  \activateIncludeBanner{#1}%%
}{%%
  \deactivateIncludeBanner%%
}

\makeatletter
\apptocmd{\@chapter}{\ifincludeBannerActive\thispagestyle{include-banner}\fi}{}{}
\apptocmd{\@schapter}{\ifincludeBannerActive\thispagestyle{include-banner}\fi}{}{}
\makeatother



% --------------------------------------------------
% Part title pages: vertically offset, left-aligned, no header/footer
% --------------------------------------------------
\usepackage{ragged2e} % make sure this is present near the top

\newenvironment{parttitlepage}{%
  \clearpage
  \thispagestyle{empty}% no running head / footer
  \markboth{}{}% clear left/right marks
  \vspace*{0.32\textheight}% push content down a bit
  \begingroup
  \RaggedRight% left-aligned, no justification
}{%
  \par
  \endgroup
  \clearpage
}



% --------------------------------------------------
% Section / part / chapter formatting
% --------------------------------------------------
\usepackage{titlesec}

% PARTS (your outer "Series"/major divisions)
% Shown on their own page, big and centered
\titleformat{\part}
  [display]
  {\centering\HeadingFont\Huge\bfseries}
  {\partname~\thepart}
  {1ex}
  {}

\titlespacing*{\part}
  {0pt}   % left
  {*6}    % before
  {*4}    % after

% CHAPTERS (your main units inside parts)
% Adjusted for a clean, book-like start
\titleformat{\chapter}
  [display]
  {\HeadingFont\LARGE\bfseries}
  {\chaptername~\thechapter}
  {1ex}
  {}

\titlespacing*{\chapter}
  {0pt}   % left
  {*4}    % before
  {*2}    % after

% SECTIONS (for your h2 / h3 hierarchy)
\titleformat{\section}
  {\HeadingFont\Large\bfseries}
  {\thesection}
  {0.5em}
  {}

\titleformat{\subsection}
  {\HeadingFont\large\bfseries}
  {\thesubsection}
  {0.5em}
  {}

% --------------------------------------------------
% Appendices
% --------------------------------------------------
\usepackage[toc,page]{appendix}

\pretocmd{\appendix}{\iaspinappendixmodetrue}{}{}

% When Pandoc/Lua filter switches to appendix mode
% (\appendix), chapters become "Appendix A", "Appendix B", etc.
% You can tweak the label if you like:
% \renewcommand{\appendixname}{Appendix}
% \renewcommand{\appendixtocname}{Appendices}

% --------------------------------------------------
% Optional: Sidebar environment for PDF
% (we can wire your epub:type="aside" to this later)
% --------------------------------------------------
\usepackage{mdframed}
\usepackage{xcolor}
\usepackage{colortbl}
\usepackage{booktabs}

\ifthenelse{\boolean{rendergrayscale}}{%%
  \selectcolormodel{gray}%%
}{}


% Sidebar background differs by mode; make it clearly visible.
\ifthenelse{\boolean{isscreen}}{%%
  % Screen PDF: slightly darker gray for clear on-screen contrast
  \colorlet{sidebarbg}{gray!10}%%
  \colorlet{calloutbg}{gray!12}%%
  \colorlet{includebannerbg}{gray!25}%%
}{%%
  % Print PDF: a bit lighter, but still clearly distinct
  \colorlet{sidebarbg}{gray!5}%%
  \colorlet{calloutbg}{gray!6}%%
  \colorlet{includebannerbg}{gray!18}%%
}

% Per-environment framing (kept in dedicated inputs for easy tweaking)
\definecolor{noteboxplainbg}{HTML}{FFF9ED}
\definecolor{noteboxplainline}{HTML}{E6C35A}
\definecolor{noteboxbg}{HTML}{FFF6D5}
\definecolor{noteboxline}{HTML}{D9A300}
\definecolor{noteboxbreadcrumbbg}{HTML}{EEF5FB}
\definecolor{noteboxbreadcrumbline}{HTML}{1D5671}
\definecolor{noteboxwarningbg}{HTML}{FFF1E6}
\definecolor{noteboxwarningline}{HTML}{C05621}
\definecolor{noteboxcontrastbg}{HTML}{F4EBFF}
\definecolor{noteboxcontrastline}{HTML}{7E3AF2}
\definecolor{noteboxtransitionbg}{HTML}{EAF2FF}
\definecolor{noteboxtransitionline}{HTML}{2563EB}
\definecolor{noteboxnavigationbg}{HTML}{E6F6F7}
\definecolor{noteboxnavigationline}{HTML}{0F6B75}
\definecolor{noteboxbridgebg}{HTML}{F4F0FF}
\definecolor{noteboxbridgeline}{HTML}{6F58B9}
\definecolor{noteboxusagebg}{HTML}{EEF4FA}
\definecolor{noteboxusageline}{HTML}{3B6EA8}
\definecolor{noteboxintegrationbg}{HTML}{E8F7F6}
\definecolor{noteboxintegrationline}{HTML}{0F766E}
\definecolor{noteboxresponsibilitybg}{HTML}{FFF4E6}
\definecolor{noteboxresponsibilityline}{HTML}{B45309}
\definecolor{noteboxreflexivebg}{HTML}{EFF2F7}
\definecolor{noteboxreflexiveline}{HTML}{64748B}
\definecolor{noteboxsynthesisbg}{HTML}{E7F7EF}
\definecolor{noteboxsynthesisline}{HTML}{2F855A}
\definecolor{noteboxdisclosurebg}{HTML}{F5EDFF}
\definecolor{noteboxdisclosureline}{HTML}{A855F7}
\definecolor{noteboxmisusebg}{HTML}{FFEDEA}
\definecolor{noteboxmisuseline}{HTML}{B42318}
\definecolor{noteboxmethodologybg}{HTML}{EEF0FF}
\definecolor{noteboxmethodologyline}{HTML}{5046E5}
\definecolor{noteboxcontextbg}{HTML}{F0F4F8}
\definecolor{noteboxcontextline}{HTML}{4A5568}
\definecolor{noteboxscopebg}{HTML}{FEF3E2}
\definecolor{noteboxscopeline}{HTML}{92400E}
\definecolor{noteboxchecklistbg}{HTML}{CCFBF1}
\definecolor{noteboxchecklistline}{HTML}{0D9488}
\definecolor{noteboxpreludebg}{HTML}{F9FAFB}
\definecolor{noteboxpreludeline}{HTML}{9CA3AF}

% --------------------------------------------------
% Table appearance helpers
% --------------------------------------------------
\definecolor{iaspTableStripe}{HTML}{F3F5F7}
\definecolor{iaspTableBorder}{HTML}{C9CED6}

% Global table row breathing room
\renewcommand{\arraystretch}{1.18}

% Standard compact: slightly smaller font and tighter padding
\newcommand{\iaspTableCompactStart}{%
  \begingroup\footnotesize%
  \setlength{\tabcolsep}{2.5pt}%
  \renewcommand{\arraystretch}{1.05}%
  % Enable hyphenation and word-wrap in table cells
  \setlength{\LTpre}{0.5em}%
  \setlength{\LTpost}{0.5em}%
}
\newcommand{\iaspTableCompactEnd}{\endgroup}

\newcommand{\iaspTableStripedStart}{\rowcolors{2}{iaspTableStripe}{white}}
\newcommand{\iaspTableStripedEnd}{\rowcolors{2}{}{}}

\newenvironment{iaspTableBordered}{%
  \begin{mdframed}[
    linecolor=iaspTableBorder,
    linewidth=0.6pt,
    backgroundcolor=white,
    skipabove=0.6em,
    skipbelow=0.6em,
    innerleftmargin=6pt,
    innerrightmargin=6pt,
    innertopmargin=6pt,
    innerbottommargin=6pt
  ]%
}{\end{mdframed}}
\definecolor{noteboxdiagnosticbg}{HTML}{E7F0FF}
\definecolor{noteboxdiagnosticline}{HTML}{2454C6}
\definecolor{noteboxclosingbg}{HTML}{F5F7FB}
\definecolor{noteboxclosingline}{HTML}{475467}
\definecolor{exampleboxbg}{HTML}{E8F1FF}
\definecolor{exampleboxline}{HTML}{2F80ED}
\definecolor{definitionboxbg}{HTML}{E6F7EF}
\definecolor{definitionboxline}{HTML}{1B7F4D}

\newcommand{\iaspConvertedGraySpec}{}
\newcommand{\iaspConvertColorToGray}[1]{%%
  \convertcolorspec{named}{#1}{gray}{\iaspConvertedGraySpec}%%
  \definecolor{#1}{gray}{\iaspConvertedGraySpec}%%
}

\ifthenelse{\boolean{rendergrayscale}}{%%
  \forcsvlist{\iaspConvertColorToGray}{%%
    noteboxplainbg,%%
    noteboxplainline,%%
    noteboxbg,%%
    noteboxline,%%
    noteboxbreadcrumbbg,%%
    noteboxbreadcrumbline,%%
    noteboxwarningbg,%%
    noteboxwarningline,%%
    noteboxcontrastbg,%%
    noteboxcontrastline,%%
    noteboxtransitionbg,%%
    noteboxtransitionline,%%
    noteboxnavigationbg,%%
    noteboxnavigationline,%%
    noteboxbridgebg,%%
    noteboxbridgeline,%%
    noteboxusagebg,%%
    noteboxusageline,%%
    noteboxintegrationbg,%%
    noteboxintegrationline,%%
    noteboxresponsibilitybg,%%
    noteboxresponsibilityline,%%
    noteboxreflexivebg,%%
    noteboxreflexiveline,%%
    noteboxsynthesisbg,%%
    noteboxsynthesisline,%%
    noteboxdisclosurebg,%%
    noteboxdisclosureline,%%
    noteboxmisusebg,%%
    noteboxmisuseline,%%
    noteboxmethodologybg,%%
    noteboxmethodologyline,%%
    noteboxcontextbg,%%
    noteboxcontextline,%%
    noteboxscopebg,%%
    noteboxscopeline,%%
    noteboxchecklistbg,%%
    noteboxchecklistline,%%
    noteboxdiagnosticbg,%%
    noteboxdiagnosticline,%%
    noteboxclosingbg,%%
    noteboxclosingline,%%
    exampleboxbg,%%
    exampleboxline,%%
    definitionboxbg,%%
    definitionboxline%%
  }%%
}{}

\providecommand{\notebadgebgnormal}{noteboxline!15!white}
\providecommand{\notebadgebgbreadcrumb}{noteboxbreadcrumbline!18!white}
\providecommand{\notebadgebgwarning}{noteboxwarningline!15!white}
\providecommand{\notebadgebgcontrast}{noteboxcontrastline!18!white}
\providecommand{\notebadgebgtransition}{noteboxtransitionline!18!white}
\providecommand{\notebadgebgnavigation}{noteboxnavigationline!18!white}
\providecommand{\notebadgebgbridge}{noteboxbridgeline!18!white}
\providecommand{\notebadgebgusage}{noteboxusageline!18!white}
\providecommand{\notebadgebgintegration}{noteboxintegrationline!18!white}
\providecommand{\notebadgebgresponsibility}{noteboxresponsibilityline!18!white}
\providecommand{\notebadgebgreflexive}{noteboxreflexiveline!18!white}
\providecommand{\notebadgebgsynthesis}{noteboxsynthesisline!18!white}
\providecommand{\notebadgebgdisclosure}{noteboxdisclosureline!18!white}
\providecommand{\notebadgebgmisuse}{noteboxmisuseline!15!white}
\providecommand{\notebadgebgmethodology}{noteboxmethodologyline!18!white}
\providecommand{\notebadgebgcontext}{noteboxcontextline!18!white}
\providecommand{\notebadgebgscope}{noteboxscopeline!18!white}
\providecommand{\notebadgebgchecklist}{noteboxchecklistline!18!white}
\providecommand{\notebadgebgdiagnostic}{noteboxdiagnosticline!18!white}
\providecommand{\notebadgebgclosing}{noteboxclosingline!18!white}
\providecommand{\sidebarbadgebg}{sidebarbg!70!white}
\providecommand{\boxbadgebg}{sidebarbg!70!white}

\providecommand{\IaspLatexDir}{assets/latex}
\input{\IaspLatexDir/box-sidebar.tex}
\input{\IaspLatexDir/box-sidebar-inline.tex}
\input{\IaspLatexDir/box-callout.tex}
\input{\IaspLatexDir/box-note.tex}
\input{\IaspLatexDir/box-example.tex}
\input{\IaspLatexDir/box-definition.tex}
\input{\IaspLatexDir/bridgehead.tex}

% --------------------------------------------------
% Hyperref styling (screen vs print)
% --------------------------------------------------
\usepackage[hidelinks]{hyperref}

\ifthenelse{\boolean{isscreen}}{%%
  % Screen PDF: colored links for easier navigation
  \hypersetup{
    breaklinks=true,
    colorlinks=true,
    linkcolor=blue,
    urlcolor=blue,
    citecolor=blue
  }
}{%%
  % Print PDF: avoid colored/boxed links
  \hypersetup{
    breaklinks=true,
    colorlinks=false,
    pdfborder={0 0 0}
  }
}

% --------------------------------------------------
% Epigraphs
% --------------------------------------------------
\usepackage{ragged2e}

\newenvironment{epigraph}{
  \begin{flushright}
  \begin{minipage}{0.65\textwidth}
  \itshape
}{
  \end{minipage}
  \end{flushright}
  \vspace{2em}
}

% Slightly tighter epigraph variant for inline, chapter-opening use
\newenvironment{epigraphtight}{
  \begin{flushright}
  \begin{minipage}{0.6\textwidth}
  \itshape
}{
  \end{minipage}
  \end{flushright}
  \vspace{1em}
}



% Usage (once your Lua filter or Pandoc mappings emit it):
% \begin{sidebar}
%   Sidebar text...
% \end{sidebar}
